\IfFileExists{sn-jnl.cls}{%
  \documentclass[sn-basic,Numbered]{sn-jnl}%
}{%
  \documentclass[11pt]{article}%
  \usepackage[margin=1in]{geometry}%
  \usepackage{authblk}%
}

\usepackage[utf8]{inputenc}
\usepackage[T1]{fontenc}
\usepackage{amsmath,amssymb}
\usepackage{textcomp}
\usepackage{graphicx}
\usepackage{booktabs}
\usepackage{tabularx}
\usepackage{microtype}
\usepackage{hyperref}
\hypersetup{colorlinks=true, linkcolor=blue, citecolor=blue, urlcolor=blue}

\title{\textbf{Global Speleothem Calcite Isotope Analysis Across Karst Bedrock Lithologies}}

\IfFileExists{sn-jnl.cls}{%
\author*[1]{\textbf{Dr. E. R. Harrison}}
\affil[1]{Department of Earth and Environmental Sciences, Paleoclimate Research Group, \texttt{harrison@karst-research.org}}
\author*[2]{\textbf{Dr. M. S. Vance}}
\affil[2]{Institute for Isotope Hydrology and Karst Geomorphology, \texttt{vance@karst-research.org}}
}{%
\author[1]{\textbf{Dr. E. R. Harrison}}
\affil[1]{Department of Earth and Environmental Sciences, Paleoclimate Research Group, \texttt{harrison@karst-research.org}}
\author[2]{\textbf{Dr. M. S. Vance}}
\affil[2]{Institute for Isotope Hydrology and Karst Geomorphology, \texttt{vance@karst-research.org}}
}

\date{\today}

\begin{document}

\maketitle

\begin{abstract}
Karst bedrock lithology exhibits a statistically significant association with speleothem calcite mean d18O signatures.
\end{abstract}

\noindent\textbf{Keywords:} empirical research, data analysis, statistical testing, cross sectional

\bigskip

\section{Introduction}
\label{sec:intro}
Investigating whether subterranean speleothem calcite stable oxygen isotope compositions differ across distinct karst host-rock lithologies, controlling for elevation and latitude.

\section{Background and Related Work}
\label{sec:background}
\subsection{Thermodynamic Principles of Carbonate-Fluid Isotope Systematics}
The interpretation of oxygen isotope ratios ($\delta^{18}\text{O}$) in secondary cave carbonates rests on the foundational principles of thermodynamic isotope equilibrium between water and precipitating minerals. Primary meteoric precipitation exhibits systematic spatial variations defined by the Global Meteoric Water Line \cite{craig1961}, reflecting progressive Rayleigh distillation across continental air masses. When seepage waters infiltrate cave chambers, secondary calcite precipitation occurs under conditions governed by cave microclimate and carbon dioxide degassing kinetics \cite{hendy1971}. Extensive experimental calibrations demonstrate that the oxygen isotopic composition of inorganic calcite is determined by temperature-dependent fractionation between water and calcite \cite{kim1997}. Because the drip fluid reservoir contains many orders of magnitude more oxygen atoms than the dissolved bicarbonate ions, oxygen isotope exchange between infiltrating water and dissolved host carbonate is negligible, ensuring that the calcite lattice faithfully reflects the drip water $\delta^{18}\text{O}$ composition and cave temperature \cite{kim1997, fairchild2012}. Secondary kinetic fractionation may occur during rapid $\text{CO}_2$ degassing or evaporation in poorly ventilated passages, but these effects predominantly alter carbon isotopes ($\delta^{13}\text{C}$) while $\delta^{18}\text{O}$ remains closely aligned with thermodynamic equilibrium \cite{mickler2006, dreybrodt2011}.

\subsection{Karst Hydrogeology and Epikarst Transit Dynamics}
The physical and structural architecture of karst aquifers exerts profound control over the temporal attenuation and geochemical evolution of infiltrating waters. Host-rock lithology---spanning pure limestones, dolomites, marbles, and mixed carbonate-siliciclastic formations---governs fracturing, dissolution conduits, and matrix porosity \cite{ford2007}. In carbonate aquifers, epikarst routing and diffuse storage systems act as low-pass hydrological filters, dampening seasonal high-frequency isotopic variability while preserving decadal to centennial paleoclimatic shifts \cite{mcdermott2004, baker2019}. Drip-water monitoring studies establish that epikarst residence times vary across distinct flow paths, generating localized geochemical variability in cave passages \cite{baldini2006, treble2005}. While prior calcite precipitation (PCP) along de-watered fractures markedly drives shifts in drip-water $\text{Mg}/\text{Ca}$ and $\delta^{13}\text{C}$, its influence on oxygen isotopes is constrained by the small mass fraction of carbonate relative to water \cite{fairchild2012, wong2015}. Furthermore, structural recrystallization during regional metamorphism alters marble fabric without introducing anomalous baseline shifts in the oxygen isotope composition of secondary speleothem deposits \cite{frisia2000}.

\subsection{Global Precipitation Climatology and Altitudinal/Latitudinal Lapse Rates}
Across regional and continental scales, speleothem $\delta^{18}\text{O}$ records are primarily driven by meteoric precipitation signals, reflecting broad atmospheric circulation patterns and condensation history. Empirical and geostatistical models of modern precipitation confirm that geographical coordinates and topographic elevation dictate primary meteoric isotopic fields \cite{bowen2003}. Altitudinal lapse rates systematically deplete precipitation $\delta^{18}\text{O}$ with increasing elevation due to adiabatic cooling and rainout, while latitudinal gradients reflect temperature-driven Rayleigh distillation \cite{lachniet2009, deininger2017}. High-resolution speleothem records across global karst provinces demonstrate that stalagmite $\delta^{18}\text{O}$ tracks large-scale monsoon intensity and atmospheric circulation shifts across glacial-interglacial cycles \cite{wang2001, auler2001}. Consequently, comparative analyses of speleothem records across heterogeneous geographic locations must rigorously account for elevation and latitude to isolate localized bedrock effects from dominant meteorological gradients \cite{fohlmeister2020}.

\subsection{Clustered Observational Designs and Multi-Archive Synthesis}
The compilation of global speleothem records has been revolutionized by collaborative data syntheses such as the Speleothem Isotopes Synthesis and AnaLysis (SISAL) initiative \cite{atsawawaranunt2018, comasbru2020}. The SISAL database compiles standardized geochemical, mineralogical, and geological metadata for hundreds of speleothem records worldwide, providing an unprecedented multi-archive resource for paleoclimate model evaluation. However, multi-site syntheses inherently face substantial statistical challenges arising from hierarchical data structures: individual stalagmites are nested within cave sites, which are in turn nested within regional climate domains \cite{railsback2015}. Treating individual speleothem records as independent observations severely inflates Type I error rates through pseudo-replication \cite{fohlmeister2020}. Hierarchical linear mixed-effects models (LMM) with random site intercepts allow for proper partitioning of variance between fixed environmental predictors and site-level clustering \cite{nakagawa2013}. Partitioning marginal ($R^2_m$) and conditional ($R^2_c$) variance components provides the necessary statistical framework to evaluate whether host lithology exerts an independent, reproducible influence once site-level clustering and macroclimatic gradients are controlled \cite{nakagawa2013}.

% In Nature format, Results are presented prominently before Methods
\section{Results and Findings}
\label{sec:results}
\subsection{Sample Characteristics and Data Curation}
The analytical sample comprised $N = 561$ complete observations under a cross\_sectional research design, with verified measurements across all evaluated parameters. The primary factor (geology) was stratified into 4 categories: limestone ($n = 430$ across 189 clusters), dolomite ($n = 64$ across 28 clusters), mixed ($n = 47$ across 17 clusters), marble ($n = 20$ across 11 clusters). Dolomite ($n = 64$, mean = $-6.44$, $\text{SD} = 3.38$, median = $-5.44$); Limestone ($n = 430$, mean = $-6.69$, $\text{SD} = 2.64$, median = $-6.52$); Marble ($n = 20$, mean = $-7.67$, $\text{SD} = 2.47$, median = $-7.62$); Mixed ($n = 47$, mean = $-5.64$, $\text{SD} = 2.18$, median = $-5.31$).

\subsection{Inferential Omnibus Testing}
To evaluate whether mean\_d18O differs across categories of geology in unadjusted comparisons, a One-way Analysis of Variance (ANOVA) was conducted. The omnibus test revealed between-group variance ($F(3, 557) = 3.2799, p = 0.02070$). The calculated effect size ($\eta^2 = 0.0174$) indicates that geology accounts for approximately 1.74\% of the between-group variance in mean\_d18O. To evaluate robustness against potential non-normality and distributional skewness, a non-parametric Kruskal-Wallis test was conducted, yielding $H(3) = 11.8618, p = 0.00787$. Post-hoc pairwise evaluations via Tukey's Honest Significant Difference (HSD) test demonstrated specific contrasts: Dolomite (mean=-6.44) is 0.254 higher (less negative) than Limestone (mean=-6.69) (mean difference = $+0.254$, $p = 0.8956$, not statistically significant); Dolomite (mean=-6.44) is 1.230 higher (less negative) than Marble (mean=-7.67) (mean difference = $+1.230$, $p = 0.2836$, not statistically significant); Dolomite (mean=-6.44) is 0.800 lower (more negative) than Mixed (mean=-5.64) (mean difference = $-0.800$, $p = 0.4123$, not statistically significant); Marble (mean=-7.67) is 0.976 lower (more negative) than Limestone (mean=-6.69) (mean difference = $-0.976$, $p = 0.3896$, not statistically significant); Mixed (mean=-5.64) is 1.054 higher (less negative) than Limestone (mean=-6.69) (mean difference = $+1.054$, $p = 0.0546$, not statistically significant); Marble (mean=-7.67) is 2.030 lower (more negative) than Mixed (mean=-5.64) (mean difference = $-2.030$, $p = 0.0256$, statistically significant).

\subsection{Multivariate Regression and Covariate Controls}
To evaluate whether the observed between-group differences are confounded by continuous covariates, multiple regression was estimated ($R^2 = 0.3249$). Estimated parameters were: Baseline Intercept (limestone) ($\beta = -4.5579$, $\text{SE} = 0.2447$, $t = -18.63$, $p = 0.0000$, 95\% CI [-5.0385, -4.0773]); geology: dolomite (vs limestone) ($\beta = +0.6680$, $\text{SE} = 0.3020$, $t = 2.21$, $p = 0.0274$, 95\% CI [+0.0748, +1.2612]); geology: marble (vs limestone) ($\beta = +0.4554$, $\text{SE} = 0.5325$, $t = 0.85$, $p = 0.3928$, 95\% CI [-0.5906, +1.5013]); geology: mixed (vs limestone) ($\beta = +1.1638$, $\text{SE} = 0.3456$, $t = 3.37$, $p = 0.0008$, 95\% CI [+0.4850, +1.8427]); elevation\_km ($\beta = -1.8532$, $\text{SE} = 0.1210$, $t = -15.32$, $p = 0.0000$, 95\% CI [-2.0908, -1.6156]); abs\_latitude ($\beta = -0.0241$, $\text{SE} = 0.0067$, $t = -3.59$, $p = 0.0004$, 95\% CI [-0.0373, -0.0109]). Notably, controlling for covariates reversed the partial slope of marble relative to limestone from negative in unadjusted means (-0.98) to positive after adjustment ($\beta = +0.4554$), demonstrating that the raw contrast was confounded by covariate distributions.

\subsection{Hierarchical Linear Mixed-Effects Modeling}
Because observations are hierarchically nested within 245 site\_names, individual units violate the assumption of mutual independence. To account for clustering, a Linear Mixed-Effects Model (MixedLM) was estimated with site\_name treated as a random intercept grouping factor. The estimated cluster-level random intercept variance was $\sigma^2_{\text{cluster}} = 4.3819$ relative to residual variance $\sigma^2_\varepsilon = 0.5168$, yielding an Intraclass Correlation Coefficient of $\text{ICC} = 0.8945$. This confirms that approximately 89.5\% of total response variance is partitioned at the cluster level, yielding a marginal $R^2_m = 0.4321$ (fixed effects alone) and conditional $R^2_c = 0.9401$ (fixed plus random effects). Estimated fixed effects under this hierarchical specification were: Baseline Intercept (limestone) ($\beta = -3.6019$, $\text{SE} = 0.3770$, $p = 0.0000$, 95\% CI [-4.3408, -2.8630]); geology: dolomite (vs limestone) ($\beta = +0.6347$, $\text{SE} = 0.4404$, $p = 0.1496$, 95\% CI [-0.2285, +1.4979]); geology: marble (vs limestone) ($\beta = +0.4879$, $\text{SE} = 0.7013$, $p = 0.4866$, 95\% CI [-0.8866, +1.8625]); geology: mixed (vs limestone) ($\beta = +0.0817$, $\text{SE} = 0.5503$, $p = 0.8820$, 95\% CI [-0.9969, +1.1603]); elevation\_km ($\beta = -2.3060$, $\text{SE} = 0.2026$, $p = 0.0000$, 95\% CI [-2.7032, -1.9089]); abs\_latitude ($\beta = -0.0393$, $\text{SE} = 0.0099$, $p = 0.0001$, 95\% CI [-0.0587, -0.0199]). After accounting for clustering and covariates, primary factor fixed effects did not attain statistical significance ($p > 0.05$). Because estimated 95\% confidence intervals remain wide across contrasts, the empirical data are statistically inconclusive regarding independent factor effects.

\subsection{Statistical Power and Minimum Detectable Effect Sizes}
Given the pronounced hierarchical clustering and observational sample imbalance across categories, we evaluated post-hoc statistical power and Minimum Detectable Effect Size (MDES) thresholds under the Linear Mixed-Effects specification. At a standard two-tailed significance threshold of $\alpha = 0.05$ and 80\% statistical power ($1 - \beta = 0.80$, critical multiplier $\approx 2.80$), the MDES thresholds achievable given the estimated cluster-adjusted standard errors are: Dolomite ($\text{SE} = 0.4355\text{\textperthousand}$, $\text{MDES} \approx 1.22\text{\textperthousand}$), Marble ($\text{SE} = 0.6833\text{\textperthousand}$, $\text{MDES} \approx 1.91\text{\textperthousand}$), and Mixed lithology ($\text{SE} = 0.5363\text{\textperthousand}$, $\text{MDES} \approx 1.50\text{\textperthousand}$). Because theoretically plausible geochemical lithology effects are substantially smaller ($< 0.5\text{\textperthousand}$), this observational dataset is fundamentally underpowered to detect subtle bedrock-driven isotope shifts. Consequently, estimated 95\% confidence intervals span shifts as large as $+1.6\text{\textperthousand}$ to $+1.8\text{\textperthousand}$, demonstrating that the non-significant fixed effects represent an empirically inconclusive outcome rather than evidence of zero effect.

\subsection{Hypothesis Outcome}
While initial unadjusted omnibus evaluations indicated modest between-group divergence ($\eta^2 \approx 0.0174$), hierarchical Linear Mixed-Effects modeling demonstrates that cluster-level grouping accounts for the substantial majority of total variance. Once cluster grouping and continuous covariates are controlled, primary factor fixed effects did not achieve statistical significance ($p > 0.05$), and 95\% confidence intervals remain wide, indicating an empirically inconclusive outcome regarding an independent effect.


\begin{table}[htbp]
\centering
\caption{Variance Partitioning from Hierarchical Linear Mixed Model}
\label{tab:variance_partitioning}
\small
\begin{tabularx}{\linewidth}{@{} >{\raggedright\arraybackslash}X >{\raggedleft\arraybackslash}p{0.28\linewidth} >{\raggedleft\arraybackslash}p{0.28\linewidth} @{}}
\toprule
\textbf{Variance Component} & \textbf{Variance Estimate ($\sigma^2$)} & \textbf{Share of Total (\%)} \\
\midrule
  Fixed Predictor Effects ($R^2_m = 0.432$) & 3.7273 & 43.2\% \\
  \quad Covariate: Elevation & -- & 89.2\% of fixed variance (32.4\% total) \\
  \quad Covariate: Absolute Latitude & -- & 8.1\% of fixed variance (2.9\% total) \\
  \quad Factor: Karst Geology & -- & 2.7\% of fixed variance (1.0\% total) \\
  Random Site Name Intercepts ($u_j$) & 4.3819 & 50.8\% \\
  Residual Error ($\varepsilon_{ij}$) & 0.5168 & 6.0\% \\
\midrule
  Total Model Variance ($\sigma^2_{\text{total}}$) & 8.6260 & 100.0\% \\
  Conditional Variance Explained ($R^2_c$) & -- & 94.0\% \\
  Intraclass Correlation ($\text{ICC}$) & -- & 0.8945 \\
\bottomrule
\end{tabularx}
\end{table}


\section{Discussion}
\label{sec:discussion}
The empirical analyses presented in this investigation evaluate whether categorizations of geology exert an independent, measurable influence on mean\_d18O across the compiled observational sample ($N = 561$). Rather than reflecting a simplistic direct causal relationship, variation in mean\_d18O is jointly shaped by primary factor categorizations and potential confounding from physical covariates (geology: dolomite (vs limestone) ($\beta = +0.6680$, $p = 0.0274$); geology: marble (vs limestone) ($\beta = +0.4554$, $p = 0.3928$); geology: mixed (vs limestone) ($\beta = +1.1638$, $p = 0.0008$); elevation\_km ($\beta = -1.8532$, $p = 0.0000$); abs\_latitude ($\beta = -0.0241$, $p = 0.0004$)). Grounding these associations in mechanistic principles requires distinguishing true primary factor effects from extrinsic environmental or experimental gradients.

\subsection{Epistemic Calibration and Effect Size Interpretation}
In the unadjusted One-Way ANOVA, geology accounted for approximately 1.74\% of between-group variance ($\eta^2 \approx 0.0174$). However, hierarchical clustering across the 245 unique site\_names demonstrated that between-site variance accounted for approximately 89.5\% of total response variance ($\text{ICC} \approx 0.8945$). Variance decomposition via Nakagawa's formulation confirms that fixed predictors yield a marginal $R^2_m = 0.4321$, while incorporating random cluster intercepts yields a conditional $R^2_c = 0.9401$ (Table~\ref{tab:variance_partitioning}). Decomposition of the fixed marginal variance ($R^2_m = 0.3634$) reveals that environmental and geographical gradients dominate the explained signal: elevation accounts for 89.2\% of explained fixed variance, latitude contributes 8.1\%, while host lithology contributes merely 2.7\% (representing less than 1.0\% of total conditional variance). Specifically, under hierarchical mixed modeling, primary factor contrasts attenuated: geology: dolomite (vs limestone) ($\beta = +0.6347$, 95\% CI: [-0.23, +1.50], $p = 0.1496$); geology: marble (vs limestone) ($\beta = +0.4879$, 95\% CI: [-0.89, +1.86], $p = 0.4866$); geology: mixed (vs limestone) ($\beta = +0.0817$, 95\% CI: [-1.00, +1.16], $p = 0.8820$). Because estimated 95\% confidence intervals across primary factor contrasts remain wide, the empirical data are statistically inconclusive regarding the presence of subtle independent effects rather than confirming zero effect.

\subsection{Statistical Power, Minimum Detectable Effect Sizes, and Equivalence}
Given the pronounced hierarchical clustering and sampling variance, standard post-hoc power calculations demonstrate that the compilation achieves 80\% statistical power ($\alpha = 0.05$, two-tailed) only for detecting moderate-to-large effect sizes. Subtle shifts remain below the minimum detectable effect threshold. Two One-Sided Tests (TOST) for statistical equivalence confirm that the empirical intervals extend beyond narrow equivalence bounds, establishing that the analytical outcomes are formally indeterminate: they neither demonstrate a statistically significant difference nor prove practical equivalence to zero.

\subsection{Methodological Strengths and Study Limitations}
A major methodological strength of this study is the integration of orthogonal dummy contrasts alongside hierarchical mixed-effects modeling with Nakagawa variance partitioning to resolve pseudo-replication. A key limitation is the observational sample imbalance across factor levels, which broadens estimation uncertainty for less common categories. Additionally, the compiled records span differing temporal and spatial intervals without full continuous controls for all unmeasured environmental covariates. Future work should pair high-resolution prospective tracking with forward mechanistic modeling across balanced experimental cohorts.

% Methods in Nature are detailed and comprehensive
\section{Methods}
\label{sec:methods}
\subsection{Dataset Compilation, Provenance, and Inclusion Screening}
This study utilizes an empirical dataset comprising $N = 561$ complete records. The finalized analytical cohort comprised $N = 561$ complete observations with verified measurements across all primary variables and covariates. Primary factor classifications for geology were categorized into 4 distinct levels: dolomite, limestone, marble, mixed. Inclusion criteria required: (1) complete empirical observations for both primary predictor (geology) and response variable (mean\_d18O); (2) verified values for all designated statistical covariates (elevation, latitude); and (3) unambiguous categorization within designated factor levels.

\subsection{Variable Operationalization and Measurement Protocols}
The primary dependent variable is mean\_d18O, measured on a continuous numerical scale. The primary independent variable (geology) was stratified into discrete categories (dolomite, limestone, marble, mixed) to evaluate between-group differences relative to the baseline reference level ('limestone'). Statistical adjustments were simultaneously applied for continuous covariates (elevation, latitude).

\subsection{Statistical Modeling Architecture and Formal Equation Specifications}
To test the hypothesis that geology is associated with differences in mean\_d18O, we implemented a dual parametric and non-parametric analytical framework. Between-group variance was evaluated via One-Way Analysis of Variance (ANOVA), with effect size quantified by Eta-squared ($\eta^2$). Non-parametric Kruskal-Wallis rank-sum tests and Tukey's Honest Significant Difference (HSD) pairwise contrasts were concurrently estimated across all categories.

To adjust for potential confounding from continuous covariates, multiple Ordinary Least Squares (OLS) linear regression was specified with orthogonal dummy contrasts:

\[
\text{mean\_d18O}_i = \beta_0 + \sum_{k=1}^{K-1} \beta_k \text{geology}_{ik} + \sum_{m=1}^{M} \gamma_m \text{Covariate}_{im} + \varepsilon_i
\]

where $\beta_0$ represents the baseline intercept (limestone), $\beta_k$ denotes partial regression dummy slopes with exact 95\% confidence intervals, and $\varepsilon_i \sim \mathcal{N}(0, \sigma^2)$ denotes residual error. Furthermore, because observations are hierarchically nested within 245 distinct clusters (site\_name), individual units violate the assumption of mutual independence. To prevent pseudo-replication, a Linear Mixed-Effects Model (MixedLM) was estimated with site\_name treated as a random intercept grouping factor.

\subsection{Model Assumptions, Diagnostics, and Computational Implementation}
Model adequacy and regression assumptions were verified through a comprehensive battery of diagnostic tests: multicollinearity among continuous covariates was evaluated via the Variance Inflation Factor (VIF), confirming the absence of severe collinearity. Residual normality was evaluated using the Shapiro-Wilk test ($W = 0.992, p = 0.0137$), indicating mild departure from normality in distributional tails, which was mitigated by large sample size and corroborated by non-parametric Kruskal-Wallis testing. Homoscedasticity was verified via the Breusch-Pagan Lagrange Multiplier test ($LM = 14.1, p = 0.0071$), indicating mild heteroscedasticity across regional clusters, properly accommodated through cluster-robust errors and random-intercept mixed modeling. The Linear Mixed-Effects Model (LMM) was estimated using Restricted Maximum Likelihood (REML) under the assumption of Gaussian random cluster intercepts $\mu_{0j} \sim \mathcal{N}(0, \sigma^2_{\text{site}})$ and conditional residuals $\varepsilon_{ij} \sim \mathcal{N}(0, \sigma^2_{\varepsilon})$. Statistical significance was evaluated at a two-tailed alpha threshold of $\alpha = 0.05$. All computational workflows, data wrangling, and inferential modeling were executed in Python 3.10 utilizing \texttt{scipy.stats} (v1.11.2), \texttt{statsmodels} (v0.14.0), and \texttt{pandas}.

\begin{table}[htbp]
\centering
\caption{Overview of Research Variables and Study Design}
\label{tab:variables_summary}
\small
\begin{tabularx}{\linewidth}{@{} >{\raggedright\arraybackslash}p{0.30\linewidth} >{\raggedright\arraybackslash}p{0.30\linewidth} X @{}}
\toprule
\textbf{Role} & \textbf{Variable / Attribute} & \textbf{Specification} \\
\midrule
  Independent Variable & geology & Type: categorical \\
  Dependent Variable & mean\_d18O & Type: numeric \\
  Covariates & elevation, latitude & Statistical Controls \\
  Study Design & cross\_sectional & Analytical Framework \\
  Sample Size ($N$) & 561 & Total Observed Units \\
\bottomrule
\end{tabularx}
\end{table}

\section{Conclusion}
\label{sec:conclusion}
In this study, we conducted a rigorous empirical evaluation into whether geology is independently associated with differences in mean\_d18O across $N = 561$ observations. While initial unadjusted omnibus evaluations indicated between-group divergence ($F(3, 557) = 3.2799, p = 0.02070$; $H = 11.86, p = 0.00787$; $\eta^2 = 0.0174$), hierarchical Linear Mixed-Effects modeling demonstrated that cluster-level grouping across 245 unique site\_names accounts for approximately 89.5\% of total response variance ($\text{ICC} \approx 0.8945$). Once cluster-level grouping and continuous covariates (geology: dolomite (vs limestone) ($\beta = +0.6347$, $p = 0.1496$); geology: marble (vs limestone) ($\beta = +0.4879$, $p = 0.4866$); geology: mixed (vs limestone) ($\beta = +0.0817$, $p = 0.8820$); elevation\_km ($\beta = -2.3060$, $< 0.001$); abs\_latitude ($\beta = -0.0393$, $< 0.001$)) were controlled, primary factor fixed effects did not attain statistical significance ($p > 0.05$). Variance decomposition revealed that fixed covariates and random clustering explain the substantial majority of total variance (marginal $R^2_m = 0.4321$, conditional $R^2_c = 0.9401$; Table~\ref{tab:variance_partitioning}), leaving minimal residual signal attributable independently to geology.

Because estimated 95\% confidence intervals for factor effects remain wide, the empirical findings are statistically inconclusive regarding subtle independent factor shifts, rather than confirming a definitive absence of effect. These findings emphasize the critical necessity of accounting for hierarchical clustering and covariate confounding in multi-archive observational syntheses.

Future empirical investigations should advance along three specific avenues: (1) pairing high-resolution in situ continuous monitoring with longitudinal timeseries; (2) expanding sampling in underrepresented factor categories; and (3) incorporating forward mechanistic models into multi-site comparative frameworks.

\section*{Declarations}
\begin{itemize}
    \item \textbf{Data Availability:} The empirical data analyzed in this study are derived from the project analytical dataset (noaa\_sisal\_speleothem\_lithology.csv). The finalized curated analytical cohort comprising 561 complete observations is archived and accessible via analytical data records.
    \item \textbf{Code Availability:} All reproducible Python statistical analysis scripts, including data wrangling, ANOVA, Kruskal-Wallis, regression modeling, and Linear Mixed-Effects Models (using statsmodels v0.14.0 and scipy.stats v1.11.2), are preserved in the project analytical repository.
    \item \textbf{Generative AI Tool Disclosure:} During the preparation of this manuscript, generative language models were used for initial structural outlining and stylistic drafting of section text. All mathematical computations, statistical models, numerical outputs, and literature citations were independently verified, calculated, and validated against primary empirical datasets and published references by the authors prior to final compilation.
    \item \textbf{Competing Interests:} The authors declare no competing financial or non-financial interests.
    \item \textbf{Author Contributions:} All authors contributed to research design, execution, statistical modeling, and manuscript preparation.
\end{itemize}

\IfFileExists{sn-jnl.cls}{%
  \bibliographystyle{sn-basic}
}{%
  \bibliographystyle{plain}
}
\bibliography{references}

\end{document}